\section{Approach}\label{sec:approach}

\begin{figure*}[h]
  \centering
  \includegraphics[width=0.96\linewidth]{figure/01-workflow.png}
  \caption{Overview of the MR-based behavioral validation framework.}
  \label{fig:workflow}
\end{figure*}

\subsection{Overview}\label{sec:approach:overview}
~\cref{fig:workflow} presents the overall framework. Given a UI component $C$ and its existing test suite $T(C)$, the framework first uses a UI MR taxonomy to construct a component-specific MR space $M(C)$ from the component's public API, documentation, and implementation. It then aligns the inferred MRs with existing tests to determine which behavioral relations are exercised and which are explicitly validated. Finally, it computes MR-based completeness metrics that characterize behavioral reach and validation in the test suite. These metrics provide the basis for the empirical analyses in Section~\ref{sec:result_analysis}.


\subsection{UI MR Taxonomy}\label{sec:approach:taxonomy}

To provide a structured representation of UI component behavior, we derive a UI-specific MR taxonomy by synthesizing behavioral dimensions commonly reflected in component APIs, official documentation, GUI testing, accessibility guidelines, and metamorphic testing. We consolidate related behavioral concepts into six top-level categories: input/prop behavior, state/event semantics, interaction/accessibility, visual/layout behavior, composition/context behavior, and data flow. The \textit{Foundation} phase in \cref{fig:workflow} summarizes the taxonomy.
The taxonomy constrains MR inference to predefined behavioral relation types while allowing component-specific instantiations. This design reduces arbitrary relation generation and promotes more consistent MR inference across component libraries.

To make the subsequent inference and alignment steps concrete, \cref{fig:inferexample} presents a dropdown example. It illustrates how component evidence is used to infer a placement-consistency MR and how an existing test is later aligned with that relation. In this example, the test exercises auto-adjust placement behavior but does not assert the rendered placement outcome, forming a weak-oracle case that is revisited in Phase~1 and Phase~2.

\begin{figure}[h]
  \centering
  \scalebox{1.03}{\includegraphics[width=\linewidth]{figure/06-infer_example.png}}
  \caption{Example: placement consistency MR is inferred from the dropdown component. The existing test exercises the auto-adjust behavior (\textit{Touch}) but does not assert the rendered placement, so the relation is not \textit{Cover}.}
  \label{fig:inferexample}
\end{figure}


\subsection{Phase 1: LLM-based MR Inference}\label{sec:approach:mr-construction}

For each component, we first construct a component profile by extracting information from its public API, documentation, and implementation. The profile summarizes observable behavioral information, including configurable properties, documented behavioral expectations, state transitions, event handling, rendering logic, and accessibility semantics. Conditioned on the component profile and the predefined UI MR taxonomy, an LLM instantiates component-specific metamorphic relations. The LLM is constrained by the predefined taxonomy and required to generate structured MR descriptions consisting of a relation type, expected behavioral relation, and supporting evidence.

Each inferred MR consists of (1) a source condition, (2) an input, interaction, or state transformation, (3) an expected behavioral relation, and (4) supporting evidence extracted from the component artifacts.
As shown in \cref{fig:inferexample}, the placement consistency MR is inferred from evidence in the component profile, including placement-related API options, documentation, and implementation logic. The resulting relation specifies the source condition, transformation, expected behavioral relation, and supporting evidence used in later alignment.

The inferred relations are subsequently normalized and filtered. Relations relying on undocumented behavior, unobservable outputs, or inconsistent taxonomy labels are removed, and semantically equivalent relations are merged. The resulting MR set, denoted by $M(C)$, serves as the behavioral reference space for the subsequent completeness analysis. Its validity and taxonomy consistency are evaluated in Section~\ref{sec:validation}.

\subsection{Phase 2: Test--MR Coverage Alignment}\label{sec:approach:alignment}
Given the inferred MR space $M(C)$ and the existing test suite $T(C)$, we determine how each test relates to each inferred MR. We distinguish between two levels of behavioral validation. A test \textit{touches} an MR if it exercises the behavioral condition described by that relation. A test \textit{covers} an MR only if it additionally contains an oracle that explicitly validates the expected behavioral relation. Accordingly, we define two metrics:
\[
touch(m,T) \in \{0,1\}, \qquad cover(m,T) \in \{0,1\}.
\]

By definition, every covered relation is also touched; that is,
\[
cover(m,T)=1 \Rightarrow touch(m,T)=1.
\]

We compute these metrics using a hybrid alignment procedure that combines deterministic evidence extracted from test code with LLM-based semantic analysis. Following a disjunctive fusion strategy, an MR is labeled as \textit{Touch} or \textit{Cover} whenever either source provides supporting evidence. For \textit{Cover}, we require evidence that the test explicitly verifies the expected behavioral relation. An MR is therefore not labeled as \textit{Cover} merely because the test contains an assertion; the assertion must be relation-relevant, i.e., it must check an observable outcome implied by the inferred MR rather than simply mention or execute the corresponding behavior. Because disjunctive fusion may over-approximate alignment labels, we evaluate its precision--recall trade-off in Section~\ref{sec:validation}.

The deterministic analysis parses the test suite into individual \textit{it()}/\textit{test()} blocks, where each block defines one test case in JavaScript/TypeScript testing frameworks and searches for predefined evidence patterns relevant to the inferred relation. Such evidence is derived from the MR description and relation type and includes component configurations, event triggers, user interactions, DOM queries, callback invocations, and assertion statements.
A \textit{Touch} decision requires evidence that the behavioral condition described by the MR is exercised. A \textit{Cover} decision is more restrictive and requires evidence of (1) a relation-relevant interaction or transformation and (2) an explicit assertion validating the expected behavioral relation within the same test block. Typical oracle patterns include assertions expressed through \textit{expect()}, \textit{assert()}, or equivalent matcher APIs. The LLM-based semantic analysis applies the same behavioral criteria but reasons over the test at a semantic level, allowing it to recognize indirect oracle patterns that cannot be reliably identified through deterministic matching alone. 

For reporting and validation, alignment decisions are attributed to one of two categories. A decision is considered \textit{deterministic-supported} whenever deterministic analysis identifies supporting evidence, regardless of whether the LLM reaches the same conclusion. Otherwise, the decision is categorized as \textit{semantic-only}. 
We evaluate the reliability of the resulting alignment decisions in Section~\ref{sec:validation}. 
The output of this phase is an MR--test alignment that determines which inferred behavioral relations are exercised (\textit{Touch}) and which are explicitly validated (\textit{Cover}), providing the basis for the completeness metrics introduced in the next phase.


\subsection{Phase 3: MR Coverage Analysis}\label{sec:approach:completeness}
Using the alignment results, we compute component-level MR coverage metrics over the MR space $M(C)$. \textit{Touch} measures the proportion of inferred behavioral relations exercised by the test suite:
{\small
\[
\mathit{Touch} = \frac{|\{ m \in M(C) \mid \tch(m,T) \}|}{|M(C)|}.
\] 
}

\textit{Cover} measures the proportion of relations that are both exercised and explicitly validated:
{\small
\[
\mathit{Cover} = \frac{|\{ m \in M(C) \mid \cov(m,T) \}|}{|M(C)|}
\] 
}
To characterize behavioral gaps, we further distinguish two categories. An \emph{Untouched} relation is not exercised by any test:
{\small
\[
\mathit{Untouched} = \frac{|\{ m \in M(C) \mid \neg \tch(m,T) \}|}{|M(C)|}
\]}

A \emph{Weak} relation is exercised but lacks validation of the expected behavioral relation:
{\small
\[
\mathit{Weak}=\frac{
  \bigl| \{ m \in M(C) \mid \tch(m,T) \land \neg \cov(m,T) \} \bigr|
}{|M(C)|}
\]
}

These metrics distinguish missing behavioral scenarios from insufficient behavioral validation and support the analyses reported in Section~\ref{sec:result_analysis}.

